
\documentclass{article}
\usepackage{colortbl}
\usepackage{makecell}
\usepackage{multirow}
\usepackage{supertabular}

\begin{document}

\newcounter{utterance}

\centering \large Interaction Transcript for game `eqbench', experiment `eq\_bench', episode 118 with reflexion{-}llama{-}3.1{-}8b{-}instruct.
\vspace{24pt}

{ \footnotesize  \setcounter{utterance}{1}
\setlength{\tabcolsep}{0pt}
\begin{supertabular}{c@{$\;$}|p{.15\linewidth}@{}p{.15\linewidth}p{.15\linewidth}p{.15\linewidth}p{.15\linewidth}p{.15\linewidth}}
    \# & \multicolumn{2}{c}{Player} && \multicolumn{2}{c}{Game Master} \\
    \hline

    \theutterance \stepcounter{utterance}  
    & & & \multicolumn{4}{p{0.6\linewidth}}{
        \cellcolor[rgb]{0.9,0.9,0.9}{
            \makecell[{{p{\linewidth}}}]{
                \texttt{\tiny{[P1$\langle$GM]}}
                \texttt{You have attempted this same task before and did not succeed. These are the plans you wrote after those attempts.} \\
\texttt{Trial \#1: New plan: I failed to accurately predict the emotional response of Freddy at the end of the dialogue. Upon re{-}reading the conversation, I noticed that Freddy's tone and language become increasingly defensive and dismissive as the conversation progresses. He is confronted with Izzy's harsh criticism, but instead of showing remorse or self{-}doubt, he becomes more entrenched in his views. This suggests that Freddy is likely to feel reproached, but not necessarily guilty or enlightened. Contempt is also a possible emotion, as Freddy seems to be condescending towards Izzy's views. However, the intensity of these emotions may vary. Next time, I will pay closer attention to the tone and language used by the characters to make a more accurate prediction. I will also consider the context and the power dynamics at play in the conversation.} \\
\\ 
\texttt{Trial \#2: New plan: I failed to accurately predict the emotional response of Freddy at the end of the dialogue. Upon re{-}reading the conversation, I noticed that Freddy's tone and language become increasingly defensive and dismissive as the conversation progresses. He is confronted with Izzy's harsh criticism, but instead of showing remorse or self{-}doubt, he becomes more entrenched in his views. This suggests that Freddy is likely to feel reproached, but not necessarily guilty or enlightened. Contempt is also a possible emotion, as Freddy seems to be condescending towards Izzy's views. However, the intensity of these emotions may vary.} \\
\\ 
\texttt{To improve my prediction, I will pay closer attention to the following:} \\
\\ 
\texttt{1. **Tone and language**: I will analyze the language and tone used by Freddy and Izzy to understand their emotional states. I will look for words or phrases that indicate defensiveness, condescension, or criticism.} \\
\texttt{2. **Context and power dynamics**: I will consider the context of the conversation and the power dynamics at play. In this case, Freddy is a lawyer who has a history of bending the law to his will, while Izzy is a character who is critical of his actions. This power imbalance may contribute to Freddy's defensive behavior.} \\
\texttt{3. **Character consistency**: I will look for consistency in Freddy's behavior and language throughout the conversation. If he is consistently defensive and dismissive, it is more likely that he will feel reproached or contemptuous at the end of the dialogue.} \\
\\ 
\texttt{By paying closer attention to these factors, I hope to improve my prediction of Freddy's emotional response at the end of the dialogue.} \\
\\ 
\\ 
\\ 
\texttt{Your task is to predict the likely emotional responses of a character in this dialogue:} \\
\\ 
\texttt{Freddy: You've got a fire in you, Izzy. Reminds me of a young me.} \\
\texttt{Izzy: Save your flattery, Lawson. It won't work on me.} \\
\texttt{Freddy: Who said anything about flattery? I'm merely stating an observation.} \\
\texttt{Izzy: Well, here's an observation for you. I'm not you, and I never will be.} \\
\texttt{Freddy: Oh, I'm well aware. You see, I win my cases.} \\
\texttt{Izzy: That's because you bend the law to your will. You manipulate, you lie.} \\
\texttt{Freddy: I do what it takes to win. Something you could learn from.} \\
\texttt{Izzy: I'd rather lose than become a monster like you.} \\
\texttt{Freddy: A monster? That's a bit harsh, don't you think?} \\
\texttt{Izzy: Not harsh enough. You've lost sight of what justice is.} \\
\texttt{Freddy: Justice is a luxury we can't afford, Izzy.} \\
\texttt{Izzy: No, justice is a necessity you've forgotten.} \\
\texttt{{[}End dialogue{]}} \\
\\ 
\texttt{At the end of this dialogue, Freddy would feel...} \\
\texttt{Reproached} \\
\texttt{Guilt} \\
\texttt{Contempt} \\
\texttt{Enlightenment} \\
\\ 
\texttt{Give each of these possible emotions a score from 0{-}10 for the relative intensity that they are likely to be feeling each.} \\
\\ 
\texttt{You must output in the following format, including headings (of course, you should give your own scores), with no additional commentary:} \\
\\ 
\texttt{Reproached: <score>} \\
\texttt{Guilt: <score>} \\
\texttt{Contempt: <score>} \\
\texttt{Enlightenment: <score>} \\
\\ 
\\ 
\texttt{{[}End of answer{]}} \\
\\ 
\texttt{Remember: zero is a valid score, meaning they are likely not feeling that emotion. You must score at least one emotion > 0.} \\
\\ 
\texttt{Your answer:} \\
            }
        }
    }
     \\ \\

    \theutterance \stepcounter{utterance}  
    & \multicolumn{4}{p{0.6\linewidth}}{
        \cellcolor[rgb]{0.561,0.737,0.561}{
            \makecell[{{p{\linewidth}}}]{
                \texttt{\tiny{[P1$\rangle$GM]}}
                \texttt{Reproached: 9} \\
\texttt{Guilt: 0} \\
\texttt{Contempt: 8} \\
\texttt{Enlightenment: 0} \\
            }
        }
    }
    & & \\ \\

    \theutterance \stepcounter{utterance}  
    & & & \multicolumn{2}{p{0.3\linewidth}}{
        \cellcolor[rgb]{0.9,0.9,0.9}{
            \makecell[{{p{\linewidth}}}]{
                \texttt{\tiny{[GM$|$GM]}}
                \texttt{Reproached: 9} \\
\texttt{Guilt: 0} \\
\texttt{Contempt: 8} \\
\texttt{Enlightenment: 0} \\
            }
        }
    }
    & & \\ \\

    \theutterance \stepcounter{utterance}  
    & & & \multicolumn{2}{p{0.3\linewidth}}{
        \cellcolor[rgb]{0.9,0.9,0.9}{
            \makecell[{{p{\linewidth}}}]{
                \texttt{\tiny{[GM$|$GM]}}
                \texttt{Reproached: 8} \\
\texttt{Guilt: 0} \\
\texttt{Contempt: 7} \\
\texttt{Enlightenment: 0} \\
            }
        }
    }
    & & \\ \\

    \theutterance \stepcounter{utterance}  
    & & & \multicolumn{2}{p{0.3\linewidth}}{
        \cellcolor[rgb]{0.9,0.9,0.9}{
            \makecell[{{p{\linewidth}}}]{
                \texttt{\tiny{[GM$|$GM]}}
                \texttt{game\_result = LOSE} \\
            }
        }
    }
    & & \\ \\

    \theutterance \stepcounter{utterance}  
    & & \multicolumn{4}{p{0.6\linewidth}}{
        \cellcolor[rgb]{0.95,0.95,0.95}{
            \makecell[{{p{\linewidth}}}]{
                \texttt{\tiny{[P1$\leftrightarrow$P]}}
                \texttt{Post{-}episode Reflexion (not sent to GM):} \\
\texttt{New plan: I failed to accurately predict the emotional response of Freddy at the end of the dialogue. Upon re{-}reading the conversation, I noticed that Freddy's tone and language become increasingly defensive and dismissive as the conversation progresses. He is confronted with Izzy's harsh criticism, but instead of showing remorse or self{-}doubt, he becomes more entrenched in his views. This suggests that Freddy is likely to feel reproached, but not necessarily guilty or enlightened. Contempt is also a possible emotion, as Freddy seems to be condescending towards Izzy's views. However, the intensity of these emotions may vary.} \\
\\ 
\texttt{To improve my prediction, I will pay closer attention to the following:} \\
\\ 
\texttt{1. **Tone and language**: I will analyze the language and tone used by Freddy and Izzy to understand their emotional states. I will look for words or phrases that indicate defensiveness, condescension, or criticism.} \\
\texttt{2. **Context and power dynamics**: I will consider the context of the conversation and the power dynamics at play. In this case, Freddy is a lawyer who has a history of bending the law to his will, while Izzy is a character who is critical of his actions. This power imbalance may contribute to Freddy's defensive behavior.} \\
\texttt{3. **Character consistency**: I will look for consistency in Freddy's behavior and language throughout the conversation. If he is consistently defensive and dismissive, it is more likely that he will feel reproached or contemptuous at the end of the dialogue.} \\
\\ 
\texttt{To make a more accurate prediction, I will also consider the following:} \\
\\ 
\texttt{* I will look for specific words or phrases that indicate Freddy's emotional state, such as "I'm well aware" (which suggests a defensive tone) or "A monster? That's a bit harsh, don't you think?" (which suggests condescension).} \\
\texttt{* I will analyze the conversation from Izzy's perspective to understand her emotional state and how it may be influencing Freddy's behavior.} \\
\texttt{* I will consider the implications of Freddy's actions and behavior throughout the conversation, and how they may be contributing to his emotional state at the end of the dialogue.} \\
\\ 
\texttt{By paying closer attention to these factors, I hope to improve my prediction of Freddy's emotional response at the end of the dialogue.} \\
            }
        }
    }
    & \\ \\

\end{supertabular}
}

\end{document}
